%%%%%%%%%%%%%%%%%%%%%%%%%%%%%%%%%%%%%%%%%%%%%%%%%%%%%%%%%%%%
%% Preamble — packages, commands, and environments
%%%%%%%%%%%%%%%%%%%%%%%%%%%%%%%%%%%%%%%%%%%%%%%%%%%%%%%%%%%%

%% --- Mathematics ---
\usepackage{amsmath,amsthm,amsfonts,amssymb}
\usepackage{mathtools}

%% --- Graphics and figures ---
\usepackage{graphicx}
\graphicspath{{figures/}}
\usepackage{tikz}
\usetikzlibrary{positioning,shapes.geometric,calc,arrows.meta,patterns}

%% --- Tables ---
\usepackage{tabularx}
\usepackage{booktabs}
\usepackage{array}

%% --- Algorithms ---
\usepackage{algorithm}
\usepackage{algpseudocode}

%% --- Code listings ---
\usepackage{listings}
\lstset{
  basicstyle=\ttfamily\small,
  keywordstyle=\bfseries\color{blue!70!black},
  commentstyle=\itshape\color{green!50!black},
  stringstyle=\color{red!60!black},
  numbers=left,
  numberstyle=\tiny\color{gray},
  numbersep=8pt,
  frame=single,
  breaklines=true,
  captionpos=b,
}

%% --- References ---
\usepackage{hyperref}
\usepackage{cleveref}

%% --- Index ---
\usepackage{makeidx}
\makeindex

%% --- Misc ---
\usepackage{pifont}
\usepackage{xcolor}

%% =====================================================
%% Custom mathematical commands
%% =====================================================

% Number sets
\newcommand{\N}{\mathbb{N}}
\newcommand{\Z}{\mathbb{Z}}
\newcommand{\Q}{\mathbb{Q}}
\newcommand{\R}{\mathbb{R}}
\newcommand{\C}{\mathbb{C}}
\newcommand{\F}{\mathbb{F}}
\newcommand{\Zq}{\mathbb{Z}_q}
\newcommand{\Rq}{R_q}

% Complexity
\newcommand{\bigO}{\mathcal{O}}
\newcommand{\polylog}{\operatorname{polylog}}
\newcommand{\poly}{\operatorname{poly}}
\newcommand{\negl}{\operatorname{negl}}

% Cryptographic notation
\newcommand{\Gen}{\mathsf{Gen}}
\newcommand{\Enc}{\mathsf{Enc}}
\newcommand{\Dec}{\mathsf{Dec}}
\newcommand{\Sign}{\mathsf{Sign}}
\newcommand{\Verify}{\mathsf{Verify}}
\newcommand{\KeyGen}{\mathsf{KeyGen}}
\newcommand{\Encaps}{\mathsf{Encaps}}
\newcommand{\Decaps}{\mathsf{Decaps}}

% Lattice notation
\newcommand{\lat}{\mathcal{L}}
\newcommand{\svp}{\textsc{SVP}}
\newcommand{\cvp}{\textsc{CVP}}
\newcommand{\sivp}{\textsc{SIVP}}
\newcommand{\gapsvp}{\textsc{GapSVP}}
\newcommand{\lwe}{\textsc{LWE}}
\newcommand{\rlwe}{\textsc{Ring-LWE}}
\newcommand{\mlwe}{\textsc{Module-LWE}}

% Probability and distributions
\newcommand{\getsr}{\xleftarrow{\$}}
\newcommand{\Pr}{\operatorname{Pr}}
\newcommand{\E}{\mathbb{E}}
\DeclareMathOperator{\Var}{Var}

% Norms
\newcommand{\norm}[1]{\left\lVert #1 \right\rVert}
\newcommand{\abs}[1]{\left\lvert #1 \right\rvert}
\newcommand{\ip}[2]{\left\langle #1, #2 \right\rangle}

% Misc
\newcommand{\xmark}{\ding{55}}
\newcommand{\cmark}{\ding{51}}

%% =====================================================
%% Theorem-like environments (svmono provides these,
%% but we configure them for consistency)
%% =====================================================

% svmono defines: definition, theorem, lemma, proposition,
% corollary, example, exercise, remark, proof
% They are numbered per chapter due to envcountchap option.
